\documentclass[11pt,a4paper]{article}
\usepackage[margin=26mm,headheight=14pt]{geometry}
\usepackage{fontspec}
\setmainfont{Latin Modern Roman}
\setsansfont{Latin Modern Sans}
\setmonofont{Latin Modern Mono}[Scale=MatchLowercase]
\usepackage{amsmath,amssymb,amsthm,mathtools,mathrsfs}
\usepackage{microtype,booktabs,longtable,array,enumitem}
\usepackage[authoryear,round]{natbib}
\usepackage{xcolor,xurl}
\usepackage[colorlinks=true,linkcolor=blue!45!black,citecolor=blue!45!black,urlcolor=blue!45!black]{hyperref}
\usepackage{fancyhdr}
\pagestyle{fancy}
\fancyhf{}
\lhead{\small SECTION\_GHOSTS}
\rhead{\small Incomplete research audit}
\cfoot{\thepage}
\setlength{\emergencystretch}{2em}
\setlength{\parskip}{4pt plus 1pt}
\setlength{\parindent}{0pt}
\setlist{itemsep=2pt,topsep=4pt}
\allowdisplaybreaks[2]
\newenvironment{result}[3]{%
  \par\addvspace{9pt}\phantomsection\label{result:#1}%
  \noindent\textbf{#1.\enspace #3}\quad\textsf{[#2]}\par\nobreak
}{\par\addvspace{3pt}}
\renewcommand{\proofname}{Proof}
\newcommand{\Q}{\mathbb Q}
\newcommand{\Z}{\mathbb Z}
\newcommand{\F}{\mathbb F}
\DeclareMathOperator{\Gal}{Gal}
\DeclareMathOperator{\supp}{supp}
\title{Ordinary completed unit equations and the full pronilpotent section tower\\[5pt]
\large SECTION\_GHOSTS: proofs, corrections, and the remaining limit problem}
\author{AI-produced by GPT-6 Astra (OpenAI, ChatGPT Pro mode)}
\date{18 November 2026\\\small Work in progress; incomplete}
\begin{document}
\maketitle
\begin{abstract}
We investigate whether an ordinary completed unit-equation pair for
\(\mathbb P^1_{\mathbb Q}\setminus\{0,1,\infty\}\) that admits a continuous
section of the full geometrically pronilpotent quotient must have finite
prime support. We have not proved this assertion or constructed a
counterexample. We establish the precise fixed-pair compactness reduction,
reconstruct the factorial-prime ordinary ghosts, and retain the strict
continuous-cohomology obstruction, including classes invisible at all
finite coefficient levels. Two exceptional-prime tests select a unique
class-two candidate for the original ghosts, reducing their class-three
problem to explicit scalar residue sequences. We distinguish this fixed-pair
compactness from the false compactness assertion for a varying ghost family.
For every odd coefficient prime not dividing the class number of its
cyclotomic field, including three and five, we prove that the
fixed-ramification second cohomology vanishes in every weight.
The remaining inductive condition is simultaneous local boundary matching.
The exact integral correction groups are computed, and a matching failure
above one fixed predecessor is witnessed by at most two support places.
At coefficient three, for a constructed support pattern,
odd matching layers are automatic and each adjacent odd--even pair is
controlled by a finite linear system. At weight four one fixed finite
Galois representation supplies its coefficients; the consistency criterion
for a finite construction prefix is exact, with both extendible and
excluded prefixes. This reduction does not establish that the comparisons
can all be satisfied for one fixed pair at every weight.
For a cofinal finite pro-three tower, we prove an exact one-step
criterion over the unchanged old ramification set. Its remaining
conditions are finitely many plus-character Frobenius equations.
A simultaneous arithmetic refinement preserves ordinarity at every
coefficient prime along any actual infinite path through those
equations. Such a path is not constructed; explicit frozen
predecessors fail the equations and cannot be repaired by adding
ordinary support primes.
At coefficient two we compute the fixed-ramification obstruction
in every weight and prove that prescribing the actual real point
class loses no full survivor. On an explicit support pattern,
each odd-weight matching step has one scalar modulo four per
central coordinate, with its full real-frame ambiguity retained;
each even-weight step has one finite torsion choice.
We construct one ordinary infinite-support pair with an integral
all-prime class-two section matching every off-diagonal support
point. It has a global class-three lift, but no higher matching
or full-tower survival is inferred.
For the unchanged original factorial family, a two-first-prime
construction produces a coefficient-five class-two section matching
every support point and a global class-three lift. Its next matching
gate is compatibility of explicit two-coordinate residues; each
individual local constraint can be made solvable.
We also identify the odd-prime Heisenberg-cover genera, repair the exponent-two
exception, and prove the relevant comparison for passage to covers.
The full-section consequence of a positive answer is stronger than the
prompt states: Stix's theorem already resolves the locally ordinary case
with rational abelian pair. Every unresolved assertion is retained as such;
the audit register distinguishes complete proofs from unfinished checks.
No priority is claimed, and no result is attributed to a human contributor.
\end{abstract}
\begingroup\small\tableofcontents\endgroup
\newpage
\section*{Conventions and status}
\addcontentsline{toc}{section}{Conventions and status}
We use the actual arithmetic action on the geometric fundamental group,
with a rational tangential base point when coordinates are needed.
Sections and cochains are continuous; conjugacy means geometric conjugacy.
An ordinary completed pair \((A,B)\) lies in
\(\varprojlim_n\mathbb Q^\times/(\mathbb Q^\times)^n\) in both coordinates
and has local representatives \((x_q,1-x_q)\) at every finite prime.
The power quotients have their discrete topology. A ghost is such a pair
with infinite prime support. A section with ordinary abelian localization
is not by definition a locally geometric section of a nonabelian quotient.
We distinguish those properties throughout.

The labels \emph{Proved}, \emph{Refuted}, and \emph{Known} apply to the exact
numbered statements, not to every task to which they contribute. A statement
labelled \emph{Open} records a question not settled in this run. The report
and result register give the task-level accounting. Earlier unpublished
AI-produced statements are not treated as authorities: the paper either
supplies the proof used or explicitly withholds certification. The supplied
unrefereed local section theorem is not used.

\section{P. Literature and the scope of established results}
\input{sections/part_P_literature}
\input{sections/part_P_additional_literature}
\input{sections/part_P_massey_scope}
\input{sections/part_P_regular_embedding_scope}
\section{A. The fixed-pair limit and continuous obstructions}
\input{sections/part_A_comparison}
\input{sections/part_A_cohomology}
\input{sections/part_A_projective_route}
\input{sections/part_A_density}
\input{sections/part_A_arithmetic_generated}
\section{B. The class-three remainder and finite exclusions}
\input{sections/part_B_cohomology}
\input{sections/part_B_arithmetic_generated}
\input{sections/part_B_tasks}
\section{C. Covers, rigid targets, and reciprocity}
\input{sections/part_C_geometry}
\section{D. Consequences for the manuscript assertions}
\input{sections/part_D_consequences}
\section{E. Independent proofs and the audit register}
\input{sections/part_E_arithmetic_generated}
\input{sections/part_E_finite_quotient}
\input{sections/part_E_fixed_depth}
\input{sections/part_E_later_failure}
\input{sections/part_E_multiple_punctures}
\input{sections/part_E_symmetry_audit}
\input{sections/part_E_positive_tower_audit}
\input{sections/part_E_arithmetic_cm_base}
\input{sections/part_E_fermat_audit}
\input{sections/part_E_gt_scope_audit}
\input{sections/part_E_exact_class_three}
\input{sections/part_E_exact_four}
\input{sections/part_E_coefficient_three}
\input{sections/part_E_regular_coefficients}
\input{sections/part_E_two_place_matching}
\input{sections/part_E_two_step}
\input{sections/part_E_genus_two}
\input{sections/part_E_elliptic_audit}
\input{sections/part_E_joint_action}
\input{sections/part_E_weight_four}
\input{sections/part_E_embedding_gate}
\input{sections/part_E_embedding_step}
\input{sections/part_E_everywhere_finite}
\input{sections/part_E_ordinary_invariant}
\input{sections/part_E_failed_predecessor}
\input{sections/part_E_finite_field_sections}
\input{sections/part_E_full_ordinary_bridge}
\input{sections/part_E_relative_approximation}
\input{sections/part_E_coefficient_two}
\input{sections/part_E_real_two}
\input{sections/part_E_dyadic_matching}
\input{sections/part_E_harmonic_audit}
\input{sections/part_E_two_inert_audit}
\input{sections/part_E_two_primary_root}
\input{sections/part_E_all_prime_class_two}
\input{sections/part_E_gt_geometric_repair}
\input{sections/part_E_gt_coefficient_bridge}
\input{sections/part_E_gt_tower_repair}
\input{sections/part_E_gt_two_point_marking}
\input{sections/part_E_gt_translation_fusion}
\input{sections/part_E_gt_involution_bridge}
\input{sections/part_E_gt_eight_point_parameter}
\input{sections/part_E_gt_cm_covariance}
\input{sections/part_E_gt_cm_character}
\input{sections/part_E_finite_field_detector}
\input{sections/part_E_gt_finite_separation}
\input{sections/part_E_gt_witness_normalization}
\input{sections/part_E_gt_strict_inclusion}
\input{sections/part_E_detector_global}
\input{sections/part_E_relative_cm_identity}
\input{sections/part_E_compact_quaternion_detector}
\input{sections/part_E_compact_quaternion_homology}
\input{sections/part_E_dyadic_presentation}
\input{sections/part_E_frobenius_splitting}
\input{sections/part_E_compact_catalog}
\input{sections/part_E_original_five_shell}
\input{sections/part_E_original_first_gate}
\input{sections/part_E_original_five_continuation}
\input{sections/part_E_first_gate_duality}
\input{sections/part_E_first_prime_spin}
\input{sections/part_E_original_five_matching}
\input{sections/part_E_two_first_primes}
\input{sections/part_E_register}
\input{sections/part_E_audit_table}
\bibliographystyle{apalike}
\bibliography{SECTION_GHOSTS_references}
\end{document}
